\chapter{ভূমিকা}\label{ch:intro}

\section{গবেষণার প্রেক্ষাপট}
দ্রুত বিক্রয়যোগ্য ভোগ্যপণ্য (FMCG) প্রতিষ্ঠানকে প্রতিদিন অসংখ্য ছোট সিদ্ধান্ত নিতে হয়: কোন পণ্য কতটা তৈরি হবে, কোন
ডিস্ট্রিবিউটর কোন চালান পাবে, কোন অঞ্চলে আরেকজন বিক্রয় প্রতিনিধি দরকার, কিংবা একটি প্রমোশন খরচের যোগ্য কি না। আলাদাভাবে
প্রতিটি সিদ্ধান্ত ছোট, কিন্তু সব মিলিয়ে এগুলোই মুনাফা নির্ধারণ করে। বেকারি পণ্যে আরেকটি সীমাবদ্ধতা আছে: পণ্যের মেয়াদ কম।
বেশি উৎপাদন দ্রুত অপচয়ে পরিণত হয়, আর কম উৎপাদন মানে খালি তাক। তাই ভবিষ্যৎ চাহিদা আগে জানার মূল্য এখানে অনেক বেশি
\citep{hyndman2021}।

ড্যান কেক বাংলাদেশের প্যাকেটজাত কেকের বাজারে একটি সুপরিচিত ব্র্যান্ড। এর পণ্যের মধ্যে আছে সিঙ্গেল-সার্ভ মাফিন ও লেয়ার-কেক
স্লাইস থেকে শুরু করে পাউন্ড কেক, সুইস রোল, কুকিজ ও প্রিমিয়াম গিফট বক্স। পণ্যগুলো জেনারেল ট্রেড (পাড়ার দোকান), মডার্ন ট্রেড
(সুপারমার্কেট), ডিস্ট্রিবিউটর ও ই-কমার্সের মাধ্যমে ভোক্তার কাছে পৌঁছায়। এই প্রকল্পের জন্য জানুয়ারি 2020 থেকে সেপ্টেম্বর 2026
পর্যন্ত \NRows বিক্রয় রেকর্ডের একটি লেনদেন-স্তরের ডেটাসেট সংকলন করা হয়েছে। ডেটাসেটটি এত বড় যে স্প্রেডশিটে বিশ্লেষণ ধীর
ও ভুলপ্রবণ হয়ে যায়, আবার এত ছোট যে ওপেন-সোর্স টুল দিয়ে একটি ল্যাপটপেই পুরোটা বিশ্লেষণ করা যায়। তাই এখানে প্রস্তাবিত
কাঠামো পরীক্ষার জন্য এটি একটি বাস্তবসম্মত ক্ষেত্র।

\section{সমস্যা বিবৃতি}
অনেক প্রতিষ্ঠানের নিয়মিত প্রতিবেদন এখনো হাতে তৈরি স্প্রেডশিটের উপর নির্ভরশীল। এ ধরনের প্রতিবেদন সাধারণত ব্যবস্থাপনার তিনটি
প্রশ্নের কেবল প্রথমটির উত্তর দেয়: \emph{কী ঘটেছে?} এগুলো খুব কমই ব্যাখ্যা করে \emph{কেন} ঘটেছে, আর প্রায় কখনোই অনুমান করে না
\emph{সামনে কী ঘটতে পারে}। প্রকল্প প্রস্তাবনায় নিচের নির্দিষ্ট দুর্বলতাগুলো চিহ্নিত করা হয়েছিল:
\begin{itemize}
  \item বড় ডেটাসেট হাতে বিশ্লেষণ করা কঠিন, এবং হাতের কাজ ধীর ও ভুলপ্রবণ;
  \item ঐতিহাসিক প্রতিবেদন পূর্বাভাস দেয় না, ফলে পরিকল্পনা হয় প্রতিক্রিয়াশীল;
  \item গুরুত্বপূর্ণ প্যাটার্ন, সম্পর্ক ও অস্বাভাবিকতা খুঁটিনাটির ভেতরে লুকিয়ে থাকে;
  \item ব্যবস্থাপকদের কাছে মূল কর্মদক্ষতা সূচকের (KPI) একটি সমন্বিত, হালনাগাদ চিত্র থাকে না।
\end{itemize}
তাই একটি সমন্বিত কাঠামো প্রয়োজন, যা ঐতিহাসিক ডেটা নির্ভরযোগ্যভাবে প্রক্রিয়া করবে, অর্থবহ প্যাটার্ন খুঁজে বের করবে, যাচাইকৃত
পূর্বাভাস তৈরি করবে এবং সেগুলো এমনভাবে উপস্থাপন করবে যাতে ব্যবস্থাপকরা পদক্ষেপ নিতে পারেন।

\section{লক্ষ্য ও উদ্দেশ্য}
এই প্রকল্পের \textbf{লক্ষ্য} হলো এমন একটি বুদ্ধিমান ডেটা অ্যানালিটিক্স ও মেশিন লার্নিং কাঠামো নকশা ও উন্নয়ন করা, যা কাঁচা
ব্যবসায়িক ডেটাকে কার্যকর অন্তর্দৃষ্টি ও পূর্বাভাসমূলক তথ্যে রূপান্তর করে উন্নত সিদ্ধান্ত গ্রহণে সহায়তা করে। নির্দিষ্ট উদ্দেশ্যগুলো হলো:
\begin{enumerate}[label=O\arabic*.]
  \item বিশ্লেষণের জন্য একটি কাঠামোবদ্ধ ব্যবসায়িক ডেটাসেট সংগ্রহ ও প্রস্তুত করা;
  \item অনুপস্থিত, দ্বৈত, অসংগত ও অস্বাভাবিক ডেটা সামলাতে একটি পুনঃব্যবহারযোগ্য প্রিপ্রসেসিং পাইপলাইন তৈরি করা;
  \item অনুসন্ধানমূলক ডেটা বিশ্লেষণ (EDA) করা এবং প্রবণতা, সম্পর্ক ও প্যাটার্ন চিহ্নিত করা;
  \item পূর্বাভাসমূলক বিশ্লেষণের জন্য মেশিন লার্নিং মডেল তৈরি করা এবং একাধিক অ্যালগরিদমের তুলনা করা;
  \item একটি ইন্টারঅ্যাকটিভ বিশ্লেষণী ড্যাশবোর্ড তৈরি করা;
  \item বিশ্লেষণের ফল থেকে ব্যবসা-কেন্দ্রিক অন্তর্দৃষ্টি ও সুপারিশ তৈরি করা;
  \item কাঠামোর নির্ভুলতা ও উপযোগিতা মূল্যায়ন করা এবং দেখানো যে এর ফলাফল কীভাবে ব্যবস্থাপনার সিদ্ধান্তে সহায়তা করে।
\end{enumerate}

\section{গবেষণা প্রশ্ন}
\begin{enumerate}[label=RQ\arabic*.]
  \item বিক্রয় প্রেক্ষাপটে ডেটা অ্যানালিটিক্স কৌশল কীভাবে ব্যবসায়িক সিদ্ধান্ত গ্রহণ উন্নত করতে পারে?
  \item ব্যবসায়িক লেনদেন ডেটা প্রস্তুতের জন্য কোন প্রিপ্রসেসিং কৌশলগুলো সবচেয়ে কার্যকর?
  \item নির্বাচিত সমস্যাগুলোর জন্য কোন মেশিন লার্নিং অ্যালগরিদম সবচেয়ে ভালো পূর্বাভাস দেয়?
  \item ঐতিহাসিক বিক্রয় ডেটায় কোন গুরুত্বপূর্ণ প্যাটার্ন ও সম্পর্ক পাওয়া যায়?
  \item ঐতিহাসিক ডেটা থেকে ভবিষ্যৎ ব্যবসায়িক ফলাফল কতটা নির্ভুলভাবে পূর্বাভাস করা যায়?
  \item ইন্টারঅ্যাকটিভ ড্যাশবোর্ডের মাধ্যমে বিশ্লেষণের ফল ব্যবস্থাপকদের কাছে কীভাবে কার্যকরভাবে উপস্থাপন করা যায়?
  \item প্রচলিত ঐতিহাসিক প্রতিবেদনের তুলনায় পূর্বাভাসমূলক অ্যানালিটিক্স কতটা উন্নতি আনে?
\end{enumerate}
\Cref{ch:discussion}-এ প্রতিটি প্রশ্নের সুস্পষ্ট উত্তর দেওয়া হয়েছে।

\section{পরিধি}
প্রস্তাবনায় একটি ব্যবসায়িক ক্ষেত্র বেছে নেওয়ার সুযোগ ছিল। এই গবেষণা \textbf{বিক্রয় অ্যানালিটিক্সে} কেন্দ্রীভূত, এবং এর ভেতরে চারটি
বিশ্লেষণী সমস্যায়: (i) প্রতিটি পণ্যের মাসিক চাহিদা পূর্বাভাস, (ii) বড় অর্ডারের পূর্বাভাস, (iii) পণ্য ও বাজার বিভাজন, এবং
(iv) অস্বাভাবিক দিন ও লেনদেন শনাক্তকরণ। ফলাফল একটি কার্যকর প্রোটোটাইপ, পূর্ণাঙ্গ এন্টারপ্রাইজ প্রোডাকশন সিস্টেম নয়। তবে
প্রোটোটাইপটি এমন প্রকৌশল চর্চা অনুসরণ করে যা প্রোডাকশনেও প্রযোজ্য: একটি কনফিগারেশন ফাইল, পুরো পাইপলাইন চালাতে একটি
কমান্ড, স্বয়ংক্রিয় টেস্ট এবং স্বয়ংক্রিয়ভাবে তৈরি প্রতিবেদন।

\section{অবদান}
প্রধান অবদান হলো ডেটা ইঞ্জিনিয়ারিং, পরিসংখ্যান, মেশিন লার্নিং ও বিজনেস ইন্টেলিজেন্সকে একটি পুনরুৎপাদনযোগ্য পাইপলাইনে
\emph{একীভূত} করা, যেখানে ড্যাশবোর্ড, ওয়েবসাইট ও এই প্রতিবেদনের প্রতিটি চিত্র, সারণি ও সংখ্যা একটি কমান্ডে কাঁচা ডেটা থেকে
পুনরায় তৈরি হয়। নির্দিষ্টভাবে, এই প্রকল্পের অবদান:
\begin{itemize}
  \item পণ্য-স্তরের চাহিদা পূর্বাভাসের জন্য একটি leakage-মুক্ত ফিচার নকশা ও মূল্যায়ন পদ্ধতি, যাতে আছে সময়-ক্রমিক যাচাই,
        বহু-ধাপ (recursive) পরীক্ষা এবং naive বেঞ্চমার্কের সঙ্গে তুলনা;
  \item একটি বাস্তব-আকারের ব্যবসায়িক ডেটাসেটে সাতটি রিগ্রেশন ও সাতটি ক্লাসিফিকেশন অ্যালগরিদম এবং পরিসংখ্যানিক পূর্বাভাস
        মডেলের তুলনা;
  \item একটি ব্যাখ্যা-স্তর (price--volume বিভাজন, fixed-effects elasticity, permutation importance ও SHAP), যা কেবল
        নির্ভুলতা না জানিয়ে \emph{কেন} প্রশ্নের উত্তর দেয়;
  \item বিনামূল্যে প্রকাশযোগ্য একটি ড্যাশবোর্ড ও প্রকল্প ওয়েবসাইট, এবং একটি স্বয়ংক্রিয় অন্তর্দৃষ্টি-উৎপাদক যা মেট্রিককে সুপারিশে রূপান্তর করে।
\end{itemize}

\section{প্রতিবেদনের গঠন}
\Cref{ch:literature}-এ বিজনেস অ্যানালিটিক্স, ডেটা প্রস্তুতি, পূর্বাভাস, মেশিন লার্নিং ও ড্যাশবোর্ড বিষয়ক সাহিত্য পর্যালোচনা করা
হয়েছে। \Cref{ch:method}-এ গবেষণা পদ্ধতি বর্ণিত। \Cref{ch:system}-এ কাঠামোর স্থাপত্য, ডেটাসেট ও প্রিপ্রসেসিং পাইপলাইন
বর্ণনা করা হয়েছে। \Cref{ch:eda}-এ অনুসন্ধানমূলক বিশ্লেষণ উপস্থাপিত। \Cref{ch:models}-এ মেশিন লার্নিং মডেলের উন্নয়ন ও
মূল্যায়ন। \Cref{ch:dashboard}-এ ড্যাশবোর্ড ও সিদ্ধান্ত-সহায়ক স্তর। \Cref{ch:discussion}-এ ফলাফল আলোচনা, গবেষণা প্রশ্নের
উত্তর ও সীমাবদ্ধতা। \Cref{ch:conclusion}-এ উপসংহার ও ভবিষ্যৎ কাজের প্রস্তাব।
